\documentclass{article}
\usepackage{amsmath}
\usepackage{amssymb}
\usepackage{bm} % bold for \alpha

\usepackage{titling}
\setlength{\droptitle}{-3cm}

\usepackage[utf8]{inputenc}
\usepackage[hidelinks]{hyperref}






\title{Implicit-Function Primitives in ImpliSolid}
%\vspace{-2cm}
% \title{Implicit Primitives in ImpliSolid / MP5}
\author{Sosi}
\date{}

\begin{document}
\maketitle

\noindent The pre-defined implicit functions used in \texttt{ImpliSolid},
These appear in WeDesign as primitive objects (MP5/json file-format).

\noindent
% As appear in implementations (may be mathematically wrong, but reflect the code as it is)
They reflect the actual implementation, and may have flaws.
As in \texttt{js\_iteration\_2/implicit\_function/}.
The implementator, if specified, it is credited, otherwise, it is implemented by \cite{sosi:implisolid}.

\noindent
% All primitives follow the F-REP sign convention:
Sign convention: F-Rep:
$f(\mathbf{p}) > 0$ inside the solid, $f(\mathbf{p}) < 0$ outside.
%, $f(\mathbf{p}) = 0$ on the surface.

\section{Common patterns}
\subsection{Coordinate transform}
%\label{common:homogenic-transform} -> equation-based referencing is preferred than section-based numbering

Often,
and when subclass of \texttt{transformable\_implicit\_function}
before evaluation,
% First,
the input point $\mathbf{p} = (x,y,z)$ is mapped to local coordinates $(u,v,w)$
via the stored inverse transformation matrix $M^{-1}$:
\begin{equation}
\label{eq:common:homogenous-transform:pre}% equation-based referencing is preferred than section-based numbering
  (u, v, w)^\top = M^{-1}\, (x, y, z)^\top .
\end{equation}

This is followed by evaluation of the function in local coordinates $f(u,v,w)$,
which is directly returned.

\label{common:homogenic-transform:gradient-chain-rule-output}
\paragraph{The gradient} in world coordinates is recovered via the chain rule for the
linear transformation $\nabla_{\mathbf{p}} f$,
where the matrix-vector product uses the upper-left $3\times 3$ block of $M^{-1}$:
\begin{equation}
  \nabla_{\mathbf{p}} f
  \;=\;
  (M^{-1})^\top \,
  \begin{pmatrix} \partial f/\partial u \\ \partial f/\partial v \\ \partial f/\partial w \end{pmatrix} ,
\end{equation}

\subsection{Constructor}
Often,
the default (identity) constructor
assumes $M=I$, hence, $(u,v,w) = (x,y,z)$.
% In such cases, formulas are in local coordinates, e.g. $(u,v,w)$.

\pagebreak
% ============================================================
\section{Heart}
% ============================================================

\textbf{Source:}
\texttt{js\_iteration\_2/implicit\_function/heart.hpp}
\noindent Implemented by \cite{schauvier:800358f}
and corrections by \cite{m.karotsieris:9af5c56}.


\subsubsection{Implicit function}
First, it applies the coordinate transform
of Eq. \ref{eq:common:homogenous-transform:pre}.
Then, in local frame (local coordinates),
\begin{equation}
  f(u,v,w)
  \;=\;
  u^2 w^3
  + \frac{9}{200}\,v^2 w^3
  - \left(u^2 + \frac{9}{4}\,v^2 + w^2 - 1\right)^3 .
\end{equation}

or,
\begin{align}
  f(u,v,w)
  \;&=\;
   w^3 (
    u^2 + \bm\beta\,v^2 )
  - \left( A \right)^3 ,
  \\
  A(u,v,w) &\;=\; u^2 + \bm\alpha\,v^2 + w^2 - 1
  \\
  \bm\beta &\;=\; \tfrac{9}{200} ,
  \;\;
  \bm\alpha \;=\; \tfrac{9}{4} .
\end{align}

Equivalently,
the zero set (the surface) satisfies:
\begin{equation}
  \left(u^2 + \tfrac{9}{4}\,v^2 + w^2 - 1\right)^3
  \;=\;
  u^2 w^3 + \tfrac{9}{200}\,v^2 w^3 .
\end{equation}
or,
\begin{align}
  &
  \left(u^2 + \bm\alpha\,v^2 + w^2 - 1\right)^3
  \;=\;
  ( u^2 + \bm\beta\,v^2 ) w^3 ,
  %\\
  %&
  %\bm\beta \;=\; \tfrac{9}{200}
  %, \,
  %\bm\alpha \;=\; \tfrac{9}{4}
  \\
  &
  \left(A\right)^3
  \;=\;
  ( u^2 + \bm\beta\,v^2 ) w^3 ,
\end{align}

% (Taubin 1993, 1994)
% Taubin, G. "An Accurate Algorithm for Rasterizing Algebraic Curves." In Second ACM/IEEE Symposium on Solid Modeling and Applications Proceedings. 221-230, May 1993.
% Taubin, G. "An Accurate Algorithm for Rasterizing Algebraic Curves and Surfaces." IEEE Comput. Graphics Appl. 14, 14-23, 1994.

% Also some variations:
% (Nordstrand, Kuska 2004).
%  (2x^2 + 2y^2 + z^2 - 1)^3-1/(10)x^2z^3-y^2z^3=0
% via https://mathworld.wolfram.com/HeartSurface.html

% Peterson, I. "MathTrek: Algebraic Hearts." Feb. 9, 2002.
% https://www.sciencenews.org/article/algebraic-hearts

%  (2*x^2+y^2+z^2-1)^3-(1/10)*x^2*z^3-y^2*z^3 = 0
% by (?) "tore?" "mona.hjemme"?
% https://web.archive.org/web/20260424054906/http://jalape.no/math/hearttxt

This is a variant of the \textbf{Taubin (1993) heart surface}.
The classical Taubin coefficient is $\bm\beta=\tfrac{9}{80}$; here it is $\bm\beta=\tfrac{9}{200}$,
which narrows the bottom lobes of the heart.

\subsubsection{Gradient}

Define the auxiliary quantity
\begin{equation}
  A \;=\; u^2 + \tfrac{9}{4}\,v^2 + w^2 - 1 ,
\end{equation}
so that $A^2$ is the common factor in all three gradient components.

The gradient in local coordinates $(u,v,w)$ is implemented as:
\begin{align}
  \frac{\partial f}{\partial u}
    &= 2u w^3 - 6u A^2 ,
  \\[4pt]
  \frac{\partial f}{\partial v}
    &= \frac{9}{100}\,v w^3 - \frac{27}{2}\,v A^2 ,
  \\[4pt]
  \frac{\partial f}{\partial w}
    &= 3 u^2 w^2 + \frac{27}{200}\,v^2 w^2 - 6 w A^2 .
\end{align}

The gradient in world coordinates is recovered as in section \ref{common:homogenic-transform:gradient-chain-rule-output}.

\subsubsection{Implementation notes}

\begin{itemize}
  \item Parameters \texttt{a}, \texttt{b}, \texttt{c} (radius) and
        \texttt{x}, \texttt{y}, \texttt{z} (centre) are stored by the
        constructor but \emph{not used} in the formula; the shape is
        always the canonical heart, regardless of their values.
        Only the transformation matrix $M^{-1}$ scales and positions it.
  \item The matrix constructor hard-codes $a = 6$, $b = 2.5$, $c = 1$.
  \item The variable $r = a^2$ is computed in the source but never
        referenced inside the evaluation loop.
\end{itemize}


\tableofcontents


\begin{thebibliography}{9}

\bibitem[schauvier]{schauvier:800358f}
Charles Chauvier:
Heart implicit function and gradient.
\texttt{s.chauvier} \texttt{800358f}.

\bibitem[m.karotsieris]{michael.karotsieris:9af5c56}
Michael Karotsieris:
Fixes on ``heart'' formula.
\texttt{m.karotsieris} \texttt{9af5c56}.


\bibitem[Sosi]{sosi:implisolid} Sosi, Sohail ``3D'' Siadat,
main author and implementor ImpliSolid,
unless otherwise specified.

\end{thebibliography}



\end{document}
